This chapter focuses on residential smart homes in the smart space taxonomy of \Cref{intro:tab:taxonomy}, addressing the \emph{user-interaction} dimension of control introduced in \Cref{intro:sec:background}: \emph{how} do occupants direct their smart space? We present \textbf{Control in Context}~\cite{controlincontext}, a mixed-methods study of how and why smart home users navigate the spectrum from physical device interaction to proxy-based centralized control across everyday tasks, failures, and evolving home setups. We surveyed $N=60$ smart home device owners and conducted follow-up semi-structured interviews with $N=8$ participants, finding that control is a fluid, reflective process shaped by tradeoffs between convenience, reliability, and cognitive load. Smart home platforms today are designed around vendor ecosystems and commercial integration goals, not around the actual, context-dependent ways people live with and rely on their devices---leaving users to build their own workarounds, absorb switching costs, and troubleshoot failures that the system makes no effort to surface.

This chapter is organized as follows. \Cref{cic:sec:intro} introduces our research questions and motivates the study. \Cref{cic:sec:relatedWorks-smarthome} surveys related work on smart home control and user studies. \Cref{cic:sec:methods} describes our survey and interview methodology. \Cref{cic:sec:Results-all} presents our findings, covering control scheme usage, user tradeoffs, and troubleshooting workflows. \Cref{cic:sec:discussion} discusses design implications for future smart home systems. \Cref{cic:sec:limitations} covers limitations and future work. \Cref{cic:sec:conclusion} concludes.
